% GENERATED FILE. Edit canonical lessons, then run npm run generate:books.
% canonical-source-hash: fnv1a64:de9a9bfdb2c2f951
% canonical-lessons: MW-C222-adhi, MW-C222-bhijno, MW-C222-os, MW-C222-dhujno, MW-C222-dhuvar, MW-R222-review-216-222

\chapter{Wind, Wet and Cold}
\label{ch:mw-wind-wet-and-cold}
\hlchaptermodality{222}

\begin{chapteropening}
\textbf{By the end of this chapter:} \emph{I can talk about dust storms, getting wet, dew and fog.}

Complete the last lesson of chapter 222: I can talk about dust storms, getting wet, dew and fog.
\end{chapteropening}

\section[\texorpdfstring{\mw{आंधी} (ā̃dhī)}{आंधी (ā̃dhī)}]{\mw{आंधी} (ā̃dhī) — a dust storm}

\label{lesson:MW-C222-adhi}



\begin{quote}
Before the new one: say the Marwadi for to thunder, then the Marwadi for a spell of steady rain.
\end{quote}

\subsection*{You'll want to know: \mw{आंधी}}
\textbf{\mw{आंधी}} — \emph{ā̃dhī} — \textquotedblleft{}a dust storm\textquotedblright{}.

\begin{quote}
\mw{जेठ} \mw{मांय} \mw{आंधी} \mw{आवै} \mw{है।}
\end{quote}

\begin{quote}
\emph{jeṭh mā̃y ā̃dhī āvai hai.}
\end{quote}

\begin{quote}
Dust storms come in early summer.
\end{quote}

A feminine noun.

\begin{cousinweb}
\textbf{\mw{आंधी}} (\emph{ā̃dhī}) is inherited from Sanskrit \emph{andhikā}, \textquotedblleft{}night\textquotedblright{}, perhaps for the darkness a dust storm brings; CDIAL lists Hindi and Gujarati \emph{ā̃dhī}, \textquotedblleft{}a dust storm\textquotedblright{}. CDIAL gives no Marwari form in this entry (CDIAL 390).
\end{cousinweb}

\begin{grammarlens}[title={What you've built}]
One more word for saying what comes next.
\end{grammarlens}

\subsection*{Your turn}
\begin{itemize}
\raggedright
  \item \emph{Say it:} \emph{ā̃dhī}
  \item \emph{Say it:} \emph{ā̃dhī}, once more
  \item \emph{Say it:} \emph{jeṭh mā̃y ā̃dhī āvai hai.} — Dust storms come in early summer.
  \item \emph{Recall:} say the Marwadi for to thunder, then the Marwadi for a spell of steady rain, then say \emph{ā̃dhī} again
\end{itemize}

\begin{tcolorbox}[breakable,colback=teal!4,colframe=teal!35!black,title={Before you move on}]
What does \mw{आंधी} mean? (\textquotedblleft{}A dust storm\textquotedblright{}.) Say it once more.
\end{tcolorbox}
\section[\texorpdfstring{\mw{भीजणो} (bhījṇo)}{भीजणो (bhījṇo)}]{\mw{भीजणो} (bhījṇo) — to get wet}

\label{lesson:MW-C222-bhijno}



\begin{quote}
Before the new one: say the Marwadi for a spell of steady rain, then the Marwadi for a dust storm.
\end{quote}

\subsection*{You'll want to know: \mw{भीजणो}}
\textbf{\mw{भीजणो}} — \emph{bhījṇo} — \textquotedblleft{}to get wet\textquotedblright{}.

\begin{quote}
\mw{हूं} \mw{काल} \mw{बरसात} \mw{मांय} \mw{भीजसूं।}
\end{quote}

\begin{quote}
\emph{hū̃ kāl barsāt mā̃y bhījsū̃.}
\end{quote}

\begin{quote}
Tomorrow I'll get wet in the rain.
\end{quote}

Past (he …): \mw{वो} \mw{बरसात} \mw{मांय} \mw{भीज्यो}, \emph{vo barsāt mā̃y bhījyo}.

\begin{cousinweb}
\textbf{\mw{भीजणो}} (\emph{bhījṇo}) goes back to a reconstructed \emph{*bhiyajyate}, \textquotedblleft{}is smeared\textquotedblright{}; CDIAL lists Hindi \emph{bhījnā} and Gujarati \emph{bhijvũ}, \textquotedblleft{}to be wet\textquotedblright{}, under it. CDIAL gives no Marwari form in this entry (CDIAL 9502).
\end{cousinweb}

\begin{grammarlens}[title={What you've built}]
One more word for saying what comes next.
\end{grammarlens}

\subsection*{Your turn}
\begin{itemize}
\raggedright
  \item \emph{Say it:} \emph{bhījṇo}
  \item \emph{Say it:} \emph{bhījṇo}, once more
  \item \emph{Say it:} \emph{hū̃ kāl barsāt mā̃y bhījsū̃.} — Tomorrow I'll get wet in the rain.
  \item \emph{Recall:} say the Marwadi for a spell of steady rain, then the Marwadi for a dust storm, then say \emph{bhījṇo} again
\end{itemize}

\begin{tcolorbox}[breakable,colback=teal!4,colframe=teal!35!black,title={Before you move on}]
What does \mw{भीजणो} mean? (\textquotedblleft{}To get wet\textquotedblright{}.) Say it once more.
\end{tcolorbox}
\section[\texorpdfstring{\mw{ओस} (os)}{ओस (os)}]{\mw{ओस} (os) — dew}

\label{lesson:MW-C222-os}



\begin{quote}
Before the new one: say the Marwadi for a dust storm, then the Marwadi for to get wet.
\end{quote}

\subsection*{You'll want to know: \mw{ओस}}
\textbf{\mw{ओस}} — \emph{os} — \textquotedblleft{}dew\textquotedblright{}.

\begin{quote}
\mw{दिनूगै} \mw{खेत} \mw{मांय} \mw{ओस} \mw{है।}
\end{quote}

\begin{quote}
\emph{dinūgai khet mā̃y os hai.}
\end{quote}

\begin{quote}
Early in the morning there's dew on the field.
\end{quote}

A feminine noun.

\begin{cousinweb}
\textbf{\mw{ओस}} (\emph{os}) is inherited from Sanskrit \emph{avaśyā}, \textquotedblleft{}hoar-frost\textquotedblright{}; CDIAL lists Hindi and Gujarati \emph{os}, \textquotedblleft{}dew\textquotedblright{}. CDIAL gives no Marwari form in this entry (CDIAL 855).
\end{cousinweb}

\begin{grammarlens}[title={What you've built}]
One more word for saying what comes next.
\end{grammarlens}

\subsection*{Your turn}
\begin{itemize}
\raggedright
  \item \emph{Say it:} \emph{os}
  \item \emph{Say it:} \emph{os}, once more
  \item \emph{Say it:} \emph{dinūgai khet mā̃y os hai.} — Early in the morning there's dew on the field.
  \item \emph{Recall:} say the Marwadi for a dust storm, then the Marwadi for to get wet, then say \emph{os} again
\end{itemize}

\begin{tcolorbox}[breakable,colback=teal!4,colframe=teal!35!black,title={Before you move on}]
What does \mw{ओस} mean? (\textquotedblleft{}Dew\textquotedblright{}.) Say it once more.
\end{tcolorbox}
\section[\texorpdfstring{\mw{धूजणो} (dhūjṇo)}{धूजणो (dhūjṇo)}]{\mw{धूजणो} (dhūjṇo) — to shiver; to tremble}

\label{lesson:MW-C222-dhujno}



\begin{quote}
Before the new one: say the Marwadi for to get wet, then the Marwadi for dew.
\end{quote}

\subsection*{You'll want to know: \mw{धूजणो}}
\textbf{\mw{धूजणो}} — \emph{dhūjṇo} — \textquotedblleft{}to shiver; to tremble\textquotedblright{}.

\begin{quote}
\mw{रात} \mw{ठंडी} \mw{है} \mw{अर} \mw{हूं} \mw{धूजूं} \mw{हूं।}
\end{quote}

\begin{quote}
\emph{rāt ṭhaṇḍī hai ar hū̃ dhūjū̃ hū̃.}
\end{quote}

\begin{quote}
The night is cold and I'm shivering.
\end{quote}

Past (he …): \mw{वो} \mw{धूज्यो}, \emph{vo dhūjyo}.

\begin{cousinweb}
\textbf{\mw{धूजणो}} (\emph{dhūjṇo}) is inherited from Sanskrit \emph{dhūyate}, \textquotedblleft{}is shaken\textquotedblright{}; CDIAL lists Gujarati \emph{dhujvũ}, \textquotedblleft{}to tremble\textquotedblright{}, under it. CDIAL gives no Marwari form in this entry (CDIAL 6863).
\end{cousinweb}

\begin{grammarlens}[title={What you've built}]
One more word for saying what comes next.
\end{grammarlens}

\subsection*{Your turn}
\begin{itemize}
\raggedright
  \item \emph{Say it:} \emph{dhūjṇo}
  \item \emph{Say it:} \emph{dhūjṇo}, once more
  \item \emph{Say it:} \emph{rāt ṭhaṇḍī hai ar hū̃ dhūjū̃ hū̃.} — The night is cold and I'm shivering.
  \item \emph{Recall:} say the Marwadi for to get wet, then the Marwadi for dew, then say \emph{dhūjṇo} again
\end{itemize}

\begin{tcolorbox}[breakable,colback=teal!4,colframe=teal!35!black,title={Before you move on}]
What does \mw{धूजणो} mean? (\textquotedblleft{}To shiver; to tremble\textquotedblright{}.) Say it once more.
\end{tcolorbox}
\section[\texorpdfstring{\mw{धूंवर} (dhū̃var)}{धूंवर (dhū̃var)}]{\mw{धूंवर} (dhū̃var) — fog; mist}

\label{lesson:MW-C222-dhuvar}



\begin{quote}
Before the new one: say the Marwadi for dew, then the Marwadi for to shiver.
\end{quote}

\subsection*{You'll want to know: \mw{धूंवर}}
\textbf{\mw{धूंवर}} — \emph{dhū̃var} — \textquotedblleft{}fog; mist\textquotedblright{}.

\begin{quote}
\mw{आज} \mw{दिनूगै} \mw{धूंवर} \mw{है।}
\end{quote}

\begin{quote}
\emph{āj dinūgai dhū̃var hai.}
\end{quote}

\begin{quote}
There's fog this morning.
\end{quote}

A feminine noun.

\begin{cousinweb}
\textbf{\mw{धूंवर}} (\emph{dhū̃var}) is inherited from Sanskrit \emph{dhūmarī}, \textquotedblleft{}mist, fog\textquotedblright{}, which CDIAL sets beside a reconstructed \emph{*dhūmara}, \textquotedblleft{}smoky\textquotedblright{}, from \emph{dhūma}, \textquotedblleft{}smoke\textquotedblright{} (CDIAL 6852); CDIAL lists Gujarati \emph{dhũvar}, \textquotedblleft{}fog\textquotedblright{}. CDIAL gives no Marwari form in this entry (CDIAL 6853).
\end{cousinweb}

\begin{grammarlens}[title={What you've built}]
One more word for saying what comes next.
\end{grammarlens}

\subsection*{Your turn}
\begin{itemize}
\raggedright
  \item \emph{Say it:} \emph{dhū̃var}
  \item \emph{Say it:} \emph{dhū̃var}, once more
  \item \emph{Say it:} \emph{āj dinūgai dhū̃var hai.} — There's fog this morning.
  \item \emph{Recall:} say the Marwadi for dew, then the Marwadi for to shiver, then say \emph{dhū̃var} again
\end{itemize}

\begin{tcolorbox}[breakable,colback=teal!4,colframe=teal!35!black,title={Before you move on}]
What does \mw{धूंवर} mean? (\textquotedblleft{}Fog; mist\textquotedblright{}.) Say it once more.
\end{tcolorbox}
\section[(dialogue)]{Looking back: from Setting Out to Wind, Wet and Cold}

\label{lesson:MW-R222-review-216-222}



\begin{quote}
Say the Marwadi for to set out and a companion.
\end{quote}

\subsection*{Your turn}
Say these aloud:

\begin{itemize}
\raggedright
  \item to set out, a companion, a kos, a traveller, a camp
  \item to cross over, to climb, a burden, to load, to halt
  \item a train, a passenger, a driver, to depart, to take to a place
  \item to come (said with respect), a gift, a week, the Teej festival, in a hurry
  \item settled, ready, to put off, to change, to find out
  \item the rains, to rain, lightning, to thunder, a spell of steady rain
  \item a dust storm, to get wet, dew, to shiver, fog
\end{itemize}

\begin{tcolorbox}[breakable,colback=teal!4,colframe=teal!35!black,title={Before you move on}]
To set out? (\textbf{\mw{निकळणो}}, \emph{nikaḷṇo}.) A companion? (\textbf{\mw{साथी}}, \emph{sāthī}.) A kos? (\textbf{\mw{कोस}}, \emph{kos}.) A traveller? (\textbf{\mw{मुसाफर}}, \emph{musāphar}.) A camp? (\textbf{\mw{डेरो}}, \emph{ḍero}.) To cross over? (\textbf{\mw{लांघणो}}, \emph{lā̃ghṇo}.) To climb? (\textbf{\mw{चढ़णो}}, \emph{chaṛhṇo}.) A burden? (\textbf{\mw{बोझ}}, \emph{bojh}.) To load? (\textbf{\mw{लादणो}}, \emph{lādṇo}.) To halt? (\textbf{\mw{ठैरणो}}, \emph{ṭhairṇo}.) A train? (\textbf{\mw{रेल}}, \emph{rel}.) A passenger? (\textbf{\mw{सवारी}}, \emph{savārī}.) A driver? (\textbf{\mw{ड्राइवर}}, \emph{ḍrāivar}.) To depart? (\textbf{\mw{छूटणो}}, \emph{chhūṭṇo}.) To take to a place? (\textbf{\mw{पुगावणो}}, \emph{pugāvṇo}.) To come (said with respect)? (\textbf{\mw{पधारणो}}, \emph{padhārṇo}.) A gift? (\textbf{\mw{भेट}}, \emph{bheṭ}.) A week? (\textbf{\mw{हफ्तो}}, \emph{haphto}.) The Teej festival? (\textbf{\mw{तीज}}, \emph{tīj}.) In a hurry? (\textbf{\mw{उतावळो}}, \emph{utāvḷo}.) Settled? (\textbf{\mw{पक्को}}, \emph{pakko}.) Ready? (\textbf{\mw{त्यार}}, \emph{tyār}.) To put off? (\textbf{\mw{टालणो}}, \emph{ṭālṇo}.) To change? (\textbf{\mw{बदळणो}}, \emph{badaḷṇo}.) To find out? (\textbf{\mw{सोधणो}}, \emph{sodhṇo}.) The rains? (\textbf{\mw{मेह}}, \emph{meh}.) To rain? (\textbf{\mw{बरसणो}}, \emph{barasṇo}.) Lightning? (\textbf{\mw{बीजळी}}, \emph{bījḷī}.) To thunder? (\textbf{\mw{गाजणो}}, \emph{gājṇo}.) A spell of steady rain? (\textbf{\mw{झड़ी}}, \emph{jhaṛī}.) A dust storm? (\textbf{\mw{आंधी}}, \emph{ā̃dhī}.) To get wet? (\textbf{\mw{भीजणो}}, \emph{bhījṇo}.) Dew? (\textbf{\mw{ओस}}, \emph{os}.) To shiver? (\textbf{\mw{धूजणो}}, \emph{dhūjṇo}.) Fog? (\textbf{\mw{धूंवर}}, \emph{dhū̃var}.)
\end{tcolorbox}
